\chapter{Storage}
\label{ch:storage}

\cnans{Microsoft::Xna::Framework::Storage} is small --- three headers, two real classes and
one exception type --- and easy to overlook next to Graphics or Audio. It is also, per CNA's
own README, one of the namespaces reported at full presence alongside Graphics, Audio, and
Input/Touch, and it is worth its own short chapter precisely because it is the namespace
responsible for something every real game eventually needs: saving and loading user data.

\section{The XNA 4.0 fake-async pattern, once more}

\cnaclass{StorageDevice}'s own class documentation states its governing convention directly:
it ``uses the XNA 4.0 fake-async pattern: \texttt{BeginXxx} completes synchronously and the
paired \texttt{EndXxx} extracts the result.'' This is the same synchronous-but-async-shaped
pattern this book has already met repeatedly --- \cnaclass{LeaderboardReader}'s
\texttt{Begin} / \texttt{End} pairs (Chapter~\ref{ch:gamerservices}), \cnaclass{NetworkSession}'s own async
operations (Chapter~\ref{ch:networking}), \cnaclass{AvatarDescription::BeginGetFromGamer} (Chapter~\ref{ch:avatar}) --- and
\cnaclass{StorageDevice} is arguably where the pattern originated in spirit: real XNA's
storage API had to accommodate a genuinely slow, UI-driven device-selection flow on a real
Xbox 360 (a player physically choosing which memory unit or hard drive to save to), even
though CNA's own desktop-filesystem-backed implementation completes every one of these calls
instantly.

\section{StorageDevice: free space, connection state, and device selection}

\cnaclass{StorageDevice} exposes \texttt{FreeSpace} and \texttt{TotalSpace} (both in bytes,
reporting real values from the underlying filesystem) and \texttt{IsConnected}. Opening a
container follows the fake-async pattern:
\texttt{BeginOpenContainer(displayName, callback, state)} / \texttt{EndOpenContainer(result)};
\texttt{DeleteContainer(titleName)} removes an entire container's directory tree outright.
Four overloads of the static \texttt{BeginShowSelector} / \texttt{EndShowSelector} pair cover
every combination real XNA supports --- with or without a specific \cnaclass{PlayerIndex},
with or without an explicit required free-space and directory-count --- for showing the
device-selection UI, even though, again, there is only ever one real ``device'' (the local
filesystem) for this to resolve to.

Two \texttt{NOXNA} static methods exist specifically because CNA has no Xbox-360-style
concept of a physically distinct storage device to select in the first place:
\texttt{SetAppNameEXT(appName)} sets the application name used to build a real,
per-application storage root directory on the host filesystem, and
\texttt{GetStorageRootEXT()} returns that resolved root path. A game is expected to call
\texttt{SetAppNameEXT} once, early --- the class's own documentation recommends the
\cnaclass{Game} constructor --- before any other storage call, since the storage root is
lazily resolved and cached on first use.

One further detail distinguishes this namespace from the other fake-async ports this book has
covered: every \texttt{Begin} / \texttt{End} result here is a real \texttt{std::unique\_ptr}, not
a caller-owned raw pointer requiring a matching manual \texttt{delete} the way
\cnaclass{NetworkSession}'s and \cnaclass{Guide}'s own \texttt{Begin} / \texttt{End} families do
(Chapters~\ref{ch:networking} and~\ref{ch:gamerservices}). \texttt{EndOpenContainer} returns a
\texttt{std::unique\_ptr<StorageContainer>}, and \texttt{EndShowSelector} returns a
\texttt{std::unique\_ptr<StorageDevice>} --- a cleaner ownership story than most of the rest of
this book's ported XNA surface, not a deviation CNA needed to call out defensively, since it
only makes the caller's life easier.

\begin{lstlisting}
// Worked example: the full device-selection-to-open-container lifecycle,
// once per application (typically in the Game constructor).
StorageDevice::SetAppNameEXT("MyGame");

auto selectResult = StorageDevice::BeginShowSelector(nullptr, nullptr);
std::unique_ptr<StorageDevice> device = StorageDevice::EndShowSelector(selectResult.get());

auto openResult = device->BeginOpenContainer("MyGame Save Data", nullptr, nullptr);
std::unique_ptr<StorageContainer> container = device->EndOpenContainer(openResult.get());
// No delete needed anywhere above -- every intermediate result is a
// unique_ptr that cleans itself up when it goes out of scope.
\end{lstlisting}

\section{StorageContainer: a real, if modest, virtual filesystem}

\cnaclass{StorageContainer} inherits \texttt{System::Object} and \texttt{System::IDisposable},
and provides exactly the file and directory operations a save-data system needs, all relative
to the container's own root: \texttt{CreateDirectory} / \texttt{DirectoryExists} /
\texttt{DeleteDirectory}, \texttt{GetDirectoryNames()} (with a glob-pattern overload),
\texttt{CreateFile} / \texttt{FileExists} / \texttt{DeleteFile}, \texttt{GetFileNames()} (with the
same glob-pattern overload), and a three-overload \texttt{OpenFile} family covering file mode
alone, mode plus access, and mode plus access plus sharing --- each returning a real
\cnaclass{System::IO::Stream} (Chapter~\ref{ch:sharp-runtime-namespaces} covers \texttt{sharp-runtime}'s own \cnaclass{Stream}
type) for a game to read or write through directly. Every relative path
passed to any of these methods is resolved against the container's own root, which is itself
resolved against \texttt{StorageDevice::GetStorageRootEXT()} --- meaning the entire
namespace, end to end, ultimately roots itself in that one \texttt{NOXNA} extension point.

\begin{lstlisting}
// Worked example: writing then reading back a save file, continuing
// directly from the container opened above.
{
    std::unique_ptr<System::IO::Stream> writeStream =
        container->CreateFile("savegame.dat");
    std::vector<SharpRuntime::bytecs> saveData = SerializeGameState();
    writeStream->Write(saveData.data(), 0, static_cast<int>(saveData.size()));
    // writeStream goes out of scope here and closes/flushes automatically.
}

if (container->FileExists("savegame.dat"))
{
    std::unique_ptr<System::IO::Stream> readStream =
        container->OpenFile("savegame.dat", System::IO::FileMode::Open);
    std::vector<SharpRuntime::bytecs> loaded(readStream->getLengthProperty());
    readStream->Read(loaded.data(), 0, static_cast<int>(loaded.size()));
}

container->Dispose();
\end{lstlisting}

\subsection{A real, currently unguarded gap: relative-path resolution does not sandbox \texttt{..}}

Every method in the worked example above --- \texttt{CreateFile}, \texttt{FileExists},
\texttt{DeleteFile}, \texttt{CreateDirectory}, and the rest --- routes its \texttt{relative}
parameter through one shared private helper, \texttt{ResolvePath(relative)}, whose entire real
implementation is a single line: \texttt{(fs::path(storagePath\_) / relative).string()}. Reading
it directly confirms there is no check anywhere in this namespace, in that helper or any of its
callers, rejecting a relative path that climbs back out of the container's own root via
\texttt{..} segments --- only an empty-string check exists, guarding against a blank filename,
not a traversal attempt. \texttt{container->CreateFile("../../outside.txt")} resolves and
creates a file two directories above the container's own storage root exactly as
\texttt{std::filesystem} would normally compose that path, with no exception, no truncation, and
no sandboxing of any kind.

This matters in exactly the situation this namespace's own \texttt{OpenFile} family exists for:
a game trusting a filename that ultimately came from outside its own hardcoded save-slot
naming (a player-chosen save name, a value read from network or downloadable content, or any
other externally-influenced string handed to \texttt{CreateFile} / \texttt{OpenFile} without
first validating it) has no protection from this namespace against that filename escaping the
sandboxed storage root the whole \cnaclass{StorageContainer} abstraction is meant to represent.
A game whose own save-name input is always a fixed, hardcoded string (as the worked example
above uses) or is validated by the caller before ever reaching this API is unaffected in
practice; a game that passes an unvalidated, externally-influenced string directly through is
not. Neither the class's own header documentation nor its test suite currently mentions this
as a deliberate, scoped decision the way this book's other real gaps are --- it is a genuine,
unaddressed gap as of this writing, worth validating any externally-influenced filename against
before calling into this namespace, rather than relying on the namespace to do it for you.

\subsection{A second, compounding gap: zero automated test coverage for the whole namespace}

The \texttt{..}-traversal gap above is worth reading alongside a second fact this book has not
yet stated about \cnans{Storage}: it is the one ported XNA namespace in this whole project
with no automated test coverage at all. \cnaclass{CnaTests}, the project's shared GoogleTest
binary, mirrors \texttt{Microsoft::Xna::Framework}'s own namespace tree one directory per
namespace --- \texttt{Audio}, \texttt{Content}, \texttt{GamerServices}, \texttt{Graphics},
\texttt{Input}, \texttt{Media}, and \texttt{Net} each have their own test directory, and
several smaller namespaces get their own flat \texttt{XxxTests.cpp} file instead (this
project's own \texttt{TitleContainerTests.cpp} and \texttt{GameTests.cpp} are two examples
sitting right alongside where a \texttt{Storage} directory or file would be). No
\texttt{tests/Microsoft/Xna/Framework/Storage/} directory exists, no flat
\texttt{StorageDeviceTests.cpp} / \texttt{StorageContainerTests.cpp} exists, and grepping both
\texttt{tests/} and \texttt{examples/} for \texttt{StorageContainer} or \texttt{StorageDevice}
turns up nothing but two incidental, unrelated mentions in an unrelated demo and an unrelated
\cnans{GamerServices} test file --- confirmed directly against the actual test tree and CMake
test registration, not inferred from the absence of a \texttt{plan\_storage.md} (no such file
exists for this namespace either, so there is no dated task history to check against in the
first place).

This means \texttt{ResolvePath}'s missing \texttt{..}-guard was never something a test suite
had the opportunity to catch --- there is no regression test for \emph{any} behavior in this
namespace, guarded or not, so the gap's survival is not evidence anyone reviewed the traversal
question and accepted the risk; it is evidence no one has systematically exercised this
namespace's behavior at all yet. The same absence covers \texttt{StorageContainer}'s own small
hand-rolled glob matcher backing the \texttt{GetFileNames} / \texttt{GetDirectoryNames} search-pattern
overloads --- a real, from-scratch backtracking \texttt{*} / \texttt{?} pattern matcher (not a
thin wrapper over a library glob or a regex translation), whose correctness on edge cases like
adjacent \texttt{**} runs or a pattern with no \texttt{*} at all currently rests entirely on
manual code reading rather than a single asserted test case.

\section{Where the storage root actually resolves to}

\texttt{GetStorageRootEXT()}'s own real implementation, \texttt{EnsureStorageRoot()}, is worth
reading directly rather than trusting the one-line summary given above --- it is a real,
three-tier platform-resolution chain, not a single fixed path, and it is lazily evaluated and
cached (\texttt{storageRootInitialized\_}) exactly once per application-name change:

\begin{lstlisting}
// From EnsureStorageRoot(), the real resolution order:
// 1. SDL_GetPrefPath(nullptr, app) -- SDL3's own per-platform, per-app
//    preference-directory API (creates the directory as a side effect).
// 2. XDG_DATA_HOME/<app>, if SDL_GetPrefPath failed and the environment
//    variable is set and non-empty.
// 3. HOME/.local/share/<app>, if XDG_DATA_HOME is unset.
// 4. The current working directory, as a last-resort fallback if even
//    HOME is unset.
\end{lstlisting}

\texttt{SDL\_GetPrefPath} is the tier that actually resolves on every desktop platform this
project targets --- on Linux it typically returns something under
\texttt{\$XDG\_DATA\_HOME/<org>/<app>/} or \texttt{\textasciitilde/.local/share/<org>/<app>/} itself,
so tiers 2 and 3 exist specifically for the case where SDL's own call genuinely fails (a
sandboxed or unusual environment with no writable preference directory at all), not as the
common path. Tier 4, the current working directory, is a real fallback with a real practical
consequence worth stating plainly: a game whose save data ends up here because every preceding
tier failed is one \texttt{cd} away from writing (or reading) save data from a completely
different location the next time it runs --- the least robust of the four tiers by a wide
margin, and worth knowing about before assuming \texttt{GetStorageRootEXT()} always returns a
stable, session-independent path.

\texttt{SetAppNameEXT()}'s own implementation is worth a specific, reassuring note against a
plausible-sounding bug that turns out not to exist: calling it a \emph{second} time (contrary
to the class's own ``call once, early'' guidance) does not leave a stale, already-cached root
behind. The setter explicitly clears \texttt{storageRoot\_} and resets
\texttt{storageRootInitialized\_} to \texttt{false} every time it runs, forcing the next
\texttt{EnsureStorageRoot()} call to re-resolve from scratch under the new application name ---
a defensive detail that makes the ``call once'' guidance a recommendation for predictability,
not a hard requirement the implementation would silently violate if ignored.

\subsection{A related, smaller finding: \texttt{IsConnected} tolerates a root that does not exist yet}

\texttt{getIsConnectedProperty()}'s own real logic is worth naming precisely, because it is not
simply ``does \texttt{GetStorageRootEXT()}'s path exist'': if the resolved root itself has never
been written to yet (a brand-new install, before the first save), the property walks \emph{up}
the path one component at a time until it finds an ancestor directory that does exist, and
reports connected as long as any such ancestor is found. On a real desktop filesystem this is
almost always true --- even a completely fresh per-app directory sits under a real, already-existing
\texttt{\textasciitilde/.local/share/} or equivalent --- so \texttt{IsConnected} reporting
\texttt{true} is not itself proof the per-app storage directory has actually been created; only
a real \texttt{BeginOpenContainer} / \texttt{EndOpenContainer} call (which does create it, via
\cnaclass{StorageContainer}'s own constructor) confirms that.

\section{A real, previously-undocumented gap: \texttt{OpenFile}'s \texttt{fileShare} parameter has nowhere to go}

\cnaclass{StorageContainer}'s three-overload \texttt{OpenFile} family accepts a
\texttt{System::IO::FileShare} argument on its fullest overload, matching real XNA's own
signature exactly. Reading \texttt{StorageContainer.cpp}'s own implementation of that overload
shows the parameter is not merely unused --- it cannot be honored at all, structurally, by
anything currently in this codebase:

\begin{lstlisting}
std::unique_ptr<System::IO::Stream> StorageContainer::OpenFile(
    const std::string& file, System::IO::FileMode fileMode,
    System::IO::FileAccess fileAccess, System::IO::FileShare /*fileShare*/)
{
    if (file.empty())
        throw std::invalid_argument("Parameter file must contain a value.");
    return std::make_unique<System::IO::FileStream>(
        ResolvePath(file), fileMode, fileAccess);
}
\end{lstlisting}

The parameter name is commented out in its own declaration --- the same
\texttt{/*paramName*/} convention this project uses project-wide to mark a genuinely unused
parameter deliberately, rather than leaving a named-but-ignored parameter that reads as an
accidental oversight. Confirming this is not simply an oversight in \emph{this} function
specifically: \texttt{sharp-runtime}'s own \cnaclass{System::IO::FileStream} --- the concrete
type \texttt{OpenFile} constructs --- has exactly three constructors (path alone; path plus
\texttt{FileMode}; path plus \texttt{FileMode} plus \texttt{FileAccess}), read directly from
its own header, and none of them accepts a \texttt{FileShare} value at all. There is no
constructor overload \texttt{OpenFile} could forward a share mode to even if its own
implementation wanted to. The underlying \texttt{std::fstream} this class wraps has no built-in
concept of an OS-level share lock either, on the platforms this project currently targets, so
the gap is not one missing forwarding call away from being closed --- it would need real,
new file-locking logic added to \texttt{FileStream} itself first.

Practically, this means every value of \texttt{FileShare} --- \texttt{None} (request exclusive
access), \texttt{Read} (allow other readers only), \texttt{ReadWrite} (allow full concurrent
access) --- currently produces byte-identical behavior: whatever the host OS's own default
file-sharing semantics happen to be for a plain \texttt{std::fstream} open (unrestricted
concurrent access on every platform this project currently builds for). A game relying on
\texttt{FileShare::None} to prevent a second process from touching the same save file
concurrently --- a real, if uncommon, XNA-era pattern for guarding against a save-corruption
race --- gets no such protection here. As with the \texttt{..}-traversal gap earlier in this
chapter, this is a genuine, currently open gap rather than a documented, deliberate scope
decision; unlike that gap, this book is the first place this specific limitation has been
written down at all.

\section{StorageDeviceNotConnectedException}

\cnaclass{StorageDeviceNotConnectedException} is the one exception type this namespace adds,
thrown when an operation requires a connected storage device that is not currently available
--- present for API completeness with real XNA's own exception surface, even though CNA's own
always-available local-filesystem backing means this condition is rare in practice compared
to a real Xbox 360, where a memory unit could genuinely be physically removed mid-session.

\section{Why this namespace is worth remembering}

\cnans{Storage} is a good, compact illustration of a pattern this book has now shown several
times: an XNA namespace originally designed around real, physical, removable hardware and a
console's own UI conventions, faithfully reproduced in API shape, but backed by a
straightforward, always-available local filesystem underneath. The fake-async pattern, the
device-selection UI stubs, and the \texttt{NOXNA} app-name/storage-root extension together
tell the same story Chapter~\ref{ch:what-is-cna} told about the project as a whole: preserve the real API's
shape precisely, even where the platform reality underneath it is simpler than what the API
was originally built to abstract over.
